% chapters/13_rlhf_dpo.tex
% Part of: AI / LLM / Agents / Hardware Notes

\section{RLHF, DPO \& Reward Modelling}

Post-training alignment methods shape a pretrained (or SFT'd) model's behaviour to
match human preferences: being helpful, harmless, and honest. The dominant paradigms
are \textbf{Reinforcement Learning from Human Feedback (RLHF)} and its offline
successor \textbf{Direct Preference Optimization (DPO)}.

\subsection{Human preference collection}

\textbf{Why comparisons?} Absolute quality ratings (``rate this response 1--5'') are
noisy and inconsistent across annotators. Comparative judgements (``which of A or B is
better?'') are faster, more reliable, and map directly to a mathematical preference model.

\textbf{Process.}
\begin{enumerate}
  \item A prompt is sampled from a prompt distribution.
  \item Two or more responses are generated by the current model (or by different
    models / human writers).
  \item Human annotators select the preferred response, optionally with a reason.
  \item Preference pairs $(x, y_w, y_l)$ --- prompt, preferred response, rejected
    response --- are collected into a dataset.
\end{enumerate}

\textbf{Bradley-Terry model.} Most reward model training assumes the
Bradley-Terry model of pairwise preference:
\[
  P(y_w \succ y_l \mid x) = \sigma(r(x, y_w) - r(x, y_l))
\]
where $r$ is the scalar reward model score. This gives a cross-entropy loss for
reward model training.

\subsection{Reward model training}

A reward model (RM) takes a prompt-response pair $(x, y)$ and returns a scalar score
indicating how well the response aligns with human preferences.

\textbf{Architecture.} Typically initialised from the SFT checkpoint with a scalar
head replacing the language-model head. Training uses the Bradley-Terry cross-entropy
loss over preference pairs.

\textbf{Reward hacking.} Because the RM is only an approximation of true human
preference, optimising against it indefinitely leads to \emph{reward hacking}: the
policy finds responses that score highly on the RM but are not actually preferred by
humans (e.g.\ very long responses, responses with confident-sounding false statements).
A KL penalty from the SFT reference policy is applied to limit distributional drift.

\subsection{PPO for language models}

\textbf{Proximal Policy Optimization (PPO)} was used in InstructGPT (Ouyang et al.\ 2022)
to fine-tune GPT-3 with RLHF. The language model is the policy $\pi_\theta$; generating
a response is a sequence of actions; the reward is the RM score.

\textbf{Training loop:}
\begin{enumerate}
  \item Sample prompts from the training distribution.
  \item Generate responses using the current policy $\pi_\theta$ (rollout).
  \item Score responses with the frozen reward model.
  \item Apply a per-token KL penalty against the SFT reference policy $\pi_\text{SFT}$
    to the reward: $r_\text{total}(x,y) = r_\text{RM}(x,y) - \beta \cdot \text{KL}(\pi_\theta \| \pi_\text{SFT})$.
  \item Update $\pi_\theta$ using the PPO clipped objective to maximise $r_\text{total}$.
\end{enumerate}

\textbf{Challenges.} PPO is on-policy: it must generate new rollouts for every gradient
step, which is slow (decoding $\gg$ forward pass). It requires holding four large models
in memory simultaneously: policy, reference policy, reward model, and value model. This
makes PPO computationally expensive.

\subsection{Direct preference optimization (DPO)}

\begin{paperbox}[title={DPO (Rafailov et al.\ 2023)}]
  Observes that the optimal RLHF policy $\pi^*$ can be expressed analytically in terms
  of the reference policy and the reward. Substituting this relationship back into the
  reward model's Bradley-Terry loss yields a closed-form objective over preference pairs
  that does not require an explicit reward model or RL.
\end{paperbox}

\textbf{DPO loss:}
\[
  \mathcal{L}_\text{DPO} = -\mathbb{E}_{(x,y_w,y_l)}\left[
    \log \sigma\!\left(
      \beta \log \frac{\pi_\theta(y_w|x)}{\pi_\text{ref}(y_w|x)}
      - \beta \log \frac{\pi_\theta(y_l|x)}{\pi_\text{ref}(y_l|x)}
    \right)
  \right]
\]

Intuitively: increase the likelihood of the preferred response relative to the
reference, and decrease the likelihood of the rejected response relative to the
reference, with the KL constraint naturally embedded via the ratio terms.

\textbf{Advantages over PPO:}
\begin{itemize}
  \item No reward model to train separately.
  \item No on-policy rollouts; trains offline on a fixed preference dataset like SFT.
  \item Simpler codebase; fits in the same training loop as SFT.
\end{itemize}

\textbf{Limitations:} sensitive to preference dataset quality; can be unstable when
the reference policy is very different from the SFT model; does not directly improve
performance beyond the reference model's distribution (no exploration).

\subsection{GRPO \& alternatives}

\textbf{Group Relative Policy Optimization (GRPO)} was introduced by DeepSeek (Shao
et al.\ 2024) and used to train the reasoning capabilities of DeepSeek-R1.

\textbf{Key idea.} For each prompt, sample a \emph{group} of $G$ outputs from the
current policy. Compute a scalar reward $r_i$ for each (e.g.\ a verifiable correctness
signal for math). Normalise rewards within the group to get advantages:
\[
  \hat{A}_i = \frac{r_i - \text{mean}(r_{1:G})}{\text{std}(r_{1:G})}
\]
Update the policy to increase the probability of high-advantage outputs.

\textbf{Why GRPO over PPO?} PPO requires a separately-trained value model (critic) to
estimate advantages, which consumes memory and adds training complexity. GRPO replaces
the critic with group-relative normalisation: the average reward within a group of
sampled responses serves as the implicit baseline, eliminating the need for a value
model entirely. This makes GRPO roughly as memory-efficient as DPO while enabling
on-policy exploration like PPO.

\begin{notebox}
  GRPO is particularly suited to tasks with \emph{verifiable} rewards: math (check
  final answer), code (run unit tests), formal logic. For open-ended tasks where reward
  is subjective, a learned reward model must still be used.
\end{notebox}

\textbf{Other alternatives:}
\begin{itemize}
  \item \textbf{REINFORCE} (simple, high-variance baseline): GRPO is a variance-reduced
    version of REINFORCE with group baseline.
  \item \textbf{KTO} (Ethayarajh et al.\ 2024): aligns using only unpaired preference
    signals (good/bad labels without requiring $(y_w, y_l)$ pairs).
  \item \textbf{SimPO} (Meng et al.\ 2024): removes the reference model entirely;
    uses the average log-likelihood of the generated response as an implicit length
    normalisation.
\end{itemize}
